% =============================================================================
% Chapter 5: Discussion and Conclusion (3-4 pages)
% =============================================================================

% =============================================================================
% Chapter 5: Discussion and Conclusion (3-4 pages)
% =============================================================================

\section{Discussion and Conclusion}
\label{sec:conclusion}

\subsection{Summary of findings}
\label{sec:summary}

This thesis investigated the application of the QIGNN framework to
community detection through QUBO formulations of modularity
maximization. The investigation produced three main empirical
findings, summarized below.

\paragraph{The bipartition formulation matches Louvain on binary
benchmarks.} The QUBO formulation for $k=2$
(equation~\eqref{eq:qubo_k2}) reduces to a quadratic form on the
modularity matrix and contains no constraint penalties. Trained
through QIGNN, this formulation converges deterministically across
all seven binary benchmarks tested, including a graph with
$n = 1\,222$ vertices. Modularity is within $0.03$ of Louvain on
average, and normalized mutual information with the ground truth
exceeds Louvain on five of the seven benchmarks, with improvements
up to~$0.33$ on the Dolphins network.

\paragraph{The multi-class formulation has structural limitations.}
The QUBO formulation for $k > 2$ (equation~\eqref{eq:qubo_multi})
introduces $nk$ binary variables under one-hot constraints
enforced through a soft penalty. Empirically, this loss landscape
is heavily influenced by the trivial all-in-one assignment:
training collapses to it on between $10\%$ and $85\%$ of runs
across benchmarks, depending on $k$ and graph size. Orthogonality
regularization suppresses the collapse but introduces a
balance prior that conflicts with the natural structure of
imbalanced real-world graphs, reducing modularity below the
unregularized baseline on Polbooks (drop from $0.46$ to $0.06$)
and on Email-EU-core (drop from $0.08$ to $0.001$). Hyperparameter
tuning does not close the gap to Louvain on any benchmark; the
gap is structural rather than calibration-related.

\paragraph{The hierarchical algorithm recovers the multi-class
case.} Recursive application of the bipartition formulation, with
modularity gain as the stopping criterion, produces a method that
matches Louvain in modularity within $0.019$ on every benchmark
in the study, including Email-EU-core
($n = 986$, $k_{\mathrm{true}} = 42$). Taking the best of the
three stopping criteria for each graph, the algorithm matches or
exceeds Louvain in NMI on \textbf{ten of thirteen benchmarks}
(eight strict wins and two ties at perfect NMI of $1.0$); the
remaining three losses are all on LFR graphs with low mixing,
where Louvain itself already attains NMI close to $1.0$. The
largest NMI gains over Louvain are on imbalanced graphs:
$+0.236$ on Polbooks-binary, $+0.094$ on Email-EU-core, and
$+0.083$ on Karate. The algorithm is empirically deterministic,
inheriting this property from the bipartition formulation: the
standard deviation of modularity across random seeds is below
$10^{-4}$ on ten of thirteen benchmarks and at most $0.005$ on
the remaining three, two orders of magnitude lower than the
multi-class formulation. The modularity-gain criterion
automatically discovers the number of communities, recovering
Louvain's choice of $k$ on graphs where Louvain's heuristic is
itself reliable.

\medskip

Two further results support the main findings. The
memory-efficient training procedure developed in
Section~\ref{sec:method_efficient} reduces the QUBO loss memory
from $O(n^2 k^2)$ to $O(n^2 + nk)$ (a factor-of-$k^2$
improvement), lowering the cost on Email-EU-core from
approximately~$6.9$~GB to approximately~$4$~MB and making the
multi-class experiments feasible on commodity hardware. The empirical analysis of Section~\ref{sec:exp_applicability}
characterizes the regime in which the methods are competitive:
the bipartition formulation works at any tested graph size for
binary tasks, and the hierarchical algorithm works at any tested
combination of $n$ and $k$, while the direct multi-class
formulation is structurally limited and is not recommended.


\subsection{Limitations}
\label{sec:limitations}

The methods developed in this thesis have several limitations
that should temper expectations.

\paragraph{Runtime.} The QIGNN-based methods run between three
and five orders of magnitude slower than Louvain. On
Email-EU-core, Louvain completes in~$0.06$ seconds while the
hierarchical algorithm requires approximately fifteen minutes.
This is the cost of training a neural network from scratch for
each (sub)graph, and it makes the methods unsuitable for
latency-sensitive applications. The runtime is a property of the
method's architecture and is unlikely to be reduced by
implementation tuning alone; sparse graph operations and
GPU acceleration would help, but a structural runtime gap to
Louvain would remain.

\paragraph{Modularity gap on the multi-class formulation.} Even
with regularization and multistart, the direct multi-class
formulation does not match Louvain in modularity on any
benchmark. The hierarchical algorithm bypasses this limitation,
but anyone wishing to apply the multi-class formulation directly
should be aware that the optimization is structurally limited.

\paragraph{Required prior $k$.} The bipartition formulation
applies only to $k=2$. The hierarchical algorithm with
modularity-gain criterion automatically determines $k$, but the
target-$k$ variant requires $k_{\mathrm{true}}$ as input.
Determining $k$ automatically when the modularity-gain criterion
is unsuitable (for instance on graphs where the modularity-optimal
partition has fewer communities than the ground truth) is left
as future work.

\paragraph{Single graph training.} Each application of the
methods requires training a GNN from scratch. The methods do not
benefit from any transfer of learned features across graphs. A
GNN trained on a corpus of graphs and applied at inference time
would amortize the training cost across many predictions; this
direction is consistent with the supervised methods of
Tsitsulin et al.~\cite{tsitsulin2020graph}, but it requires
labeled training data, which the present unsupervised setting
deliberately avoids.

\paragraph{Modularity as the only objective.} The methods
optimize modularity, which has known
limitations~\cite{fortunato2007resolution, good2010performance}.
Alternative quality functions such as normalized cut or
description-length objectives could be incorporated by adapting
the QUBO formulation, but this is outside the scope of the
thesis.


\subsection{Future work}
\label{sec:future}

The results of this thesis suggest several directions for further
research.

\paragraph{Larger graphs.} The hierarchical algorithm has been
tested up to $n = 1\,222$. Scaling to graphs of $10^4$ vertices
or more requires sparse graph operations and GPU acceleration of
the GNN training step. The structured-loss reformulation
already removes the memory bottleneck; the runtime bottleneck is
the per-bipartition training, which can be partially addressed by
GPU support and by sharing initialization across consecutive
splits.

\paragraph{Warm starts and refinement.} Louvain finds a partition
in milliseconds. A natural hybrid is to use Louvain's output as a
warm start for the hierarchical algorithm, refining only the
splits where Louvain's modularity gain was small. This combines
Louvain's speed with the hierarchical algorithm's higher
agreement with ground truth, and is consistent with the
observation that the two methods reach different local optima of
the modularity landscape.

\paragraph{Hardware QUBO solvers.} Section~\ref{sec:rw_hardware}
mentioned that the QUBO formulations of Wang et
al.~\cite{wang2024quantum} were originally solved on Coherent
Ising Machines. The structural QUBO matrices produced by the
formulations of this thesis can be directly fed to such hardware,
allowing a comparison of neural and hardware solvers on the same
problems. This comparison would clarify the practical trade-offs
between the two paradigms in the regime where both are
applicable.


\subsection{Concluding remarks}
\label{sec:final}

The original question of this thesis was whether the QIGNN
paradigm, validated on bipartite combinatorial problems such as
MaxCut and Maximum Independent Set, could be effectively applied
to the multi-class setting of community detection. The empirical
answer is twofold. The direct multi-class formulation,
constructed by the natural extension of the bipartite QUBO,
exhibits structural limitations: trivial collapse, sensitivity
to community balance, and a modularity gap to classical
heuristics that hyperparameter tuning cannot close. The
indirect approach, in which the bipartition formulation is
applied recursively, recovers the multi-class case and produces
a method that is competitive with Louvain in modularity and
typically superior to it in agreement with ground-truth
community structure.

The connection drawn in this work between QUBO-based community
detection (previously studied for quantum hardware) and
unsupervised neural solvers establishes a path toward scalable
community detection on commodity computing infrastructure. While
the methods proposed do not displace Louvain in
latency-sensitive settings, they offer a complementary tool for
applications where the quality of the partition matters more than
the speed of its computation, and they open a productive avenue
for combining QUBO-based formulations with neural network
solvers in the broader context of combinatorial optimization on
graphs.

\newpage
